\section{Gorenstein curves and the evaluation ideal}\label{scalar}
Let $R$ be a one-dimensional Gorenstein local ring. Fix a regular parameter $x$. Every nonunit nonzerodivisor generates an $\m$-primary ideal and hence divides a power of $x$. Thus $Q=Q(R)=R[1/x]$. For any module-finite extension $R\subseteq S$, set
\[
 T=H^0_{\m}(S),\qquad U=S/T,\qquad B=U\cap Q\ \subseteq S[1/x].
\]
Here $T$ consists of the elements killed by a power of $x$, so it is an ideal and meets $R$ trivially. The ring $U$ embeds in $S[1/x]$, and $B$ is a finite birational $R$-algebra. Its conductor $(R:B)$ is the ideal of all $a\in R$ such that $aB\subseteq R$.
In the intersection defining $B$, $Q$ denotes the image of the
scalar map $Q\hookrightarrow S[1/x]$.  Finiteness follows
because $B$ is an $R$-submodule of the finite module $U$.
\begin{proposition}\label{eval}
The ideal $\{\phi(1):\phi\in\Hom_R(S,R)\}$ equals $(R:B)$. In particular, the extension splits precisely when $B=R$, and $B/R=H^0_{\m}(U/R)$.
\end{proposition}
\begin{proof}
Maps to $R$ kill $T$ by regularity of $x$. The class of $u\in U$ in $U/R$ is killed by a power of $x$ exactly when $u$ belongs to $Q$. Likewise, if $x u\in B$ then $u\in B$, so $U/B$, when nonzero, has depth one.

The point where the Gorenstein hypothesis enters is extension of
functionals across $B\subseteq U$.

For any finite module $M$ on which $x$ is regular, $\Ext^1_R(M,R)=0$. Indeed $G=R/xR$ is Artinian Gorenstein, hence self-injective \cite{BH}. A free resolution of $M$ stays exact modulo $x$, since $\operatorname{Tor}_{>0}^R(M,G)=0$. Applying $\Hom_G(-,G)$ gives $\Ext^1_R(M,G)=0$. The coefficient sequence for multiplication by $x$ makes $x$ surjective on $\Ext^1_R(M,R)$, which vanishes by Nakayama. Applying this to $U/B$ shows that every functional on $B$ extends to $U$. After inverting $x$, a functional on $B$ becomes a $Q$-linear map $Q\to Q$, hence multiplication by its value at $1$. Its possible values are precisely $(R:B)$.
\end{proof}

\subsection{Generic rank and weak normality}\label{roots}
The scalar obstruction is constrained by generic rank.  A characteristic
polynomial computed on a finite free resolution supplies the necessary
integrality relations in both equal and mixed characteristic.
\begin{proposition}\label{rank}
Assume $\pd_R S<\infty$ and $\operatorname{char} k=p>0$. Then $S_Q$ is free over $Q$ of rank $r\ge1$. For $b\in B$ one has
\[
 \binom rj b^j\in R\quad(0\le j\le r),\qquad
 p^e b\in R,\qquad b^{p^e}\in R,\qquad e=v_p(r).
\]
\end{proposition}
\begin{proof}
By Auslander--Buchsbaum, $S$ admits a free resolution $0\to F_1\to F_0\to S\to0$. Over each local factor of the Artinian ring $Q$, the localized module is free by the same formula and is nonzero by scalar injectivity. Its rank is the common integer $r=\operatorname{rank} F_0-\operatorname{rank} F_1$.

Lift $b$ to $s\in S$ and lift multiplication by $s$ to endomorphisms $f_i$ of the resolution. The series $\det(1-Zf_0)/\det(1-Zf_1)$ belongs to $R[[Z]]$. Over $Q$ the resolution splits. In an adapted basis $f_0$ is block triangular with diagonal blocks $f_1$ and $bI_r$. The series therefore equals $(1-Zb)^r$ in $Q[[Z]]$. Comparing coefficients gives the first assertion. Write $r=p^e m$ with $p\nmid m$. Since $m$ is a unit of $R$, the linear coefficient gives $p^e b\in R$. In $\mathbf F_p[Z]$ the coefficient of $Z^{p^e}$ in $(1+Z)^r=(1+Z^{p^e})^m$ is $m$, so $\binom r{p^e}$ is also a unit in $R$.
\end{proof}
For $p\nmid r$ this gives $B=R$, recovering the splitting in \eqref{perfect-trace}. 

\begin{lemma}\label{wn}
Let $R$ be one-dimensional Cohen--Macaulay local of residue characteristic $p>0$. The following conditions are equivalent.
\begin{enumerate}
\item $R$ is weakly normal.
\item $q\in Q(R)$ and $q^p\in R$ imply $q\in R$.
\item $q\in Q(R)$ and $pq,q^p\in R$ imply $q\in R$.
\end{enumerate}
\end{lemma}
\begin{proof}
For $q^p\in R$ the extension $R\subseteq R[q]$ is module-finite and is an isomorphism at minimal primes, where a regular denominator for $q$ is invertible. The closed fiber is a nonzero quotient of $k[t]/(t^p-\overline{q^p})$, which has a unique prime and purely inseparable residue field. A one-dimensional local spectrum has only its minimal primes and its closed point. Thus the extension is weakly subintegral, proving (1)$\Rightarrow$(2).

Clearly (2) implies (3).  To prove the reverse implication, start with
$q^p\in R$.  The finite module $R[q]/R$ vanishes at the minimal
primes, so it has finite length.  Since $p\in\m$, there is an
$N\ge0$ with $p^Nq\in R$.  If $N>0$, set $q'=p^{N-1}q$.
Then $pq'\in R$ and $(q')^p=p^{p(N-1)}q^p\in R$, so (3)
gives $q'\in R$.  Descending induction gives $q\in R$.

Assume (2) and choose a regular parameter $x$ of $R$. A nilpotent $z\in R$ has $z/x^n\in R$ for every $n$, by repeated root closure applied to a vanishing $p$-power. Thus $z\in\bigcap_n x^nR=0$ by Krull intersection, and $R$ is reduced.
Let $C\subset Q$ be finite weakly subintegral over $R$. It is local, say with maximal ideal $\mathfrak n$, and $C/R$ has finite length. The conductor $I=(R:C)$ contains a power of $\m$ and is an ideal of $C$, so contains a power of $\mathfrak n$ since $\sqrt{\m C}=\mathfrak n$. We claim that $I$ is radical in $C$. If $b^N\in I$, then for any $c\in C$ and $p^h\ge N$ we have $(bc)^{p^h}\in I\subseteq R$. Repeated use of (2) gives $bc\in R$ for every $c$, proving $b\in I$. If $C\ne R$, then $I$ is proper and therefore $I=\mathfrak n$. For $b\in C$ some $p$-power of its residue belongs to $k$ by pure inseparability. Lifting that residue to $R$ gives $b^{p^h}\in R+\mathfrak n=R$, so $b\in R$ after all.
\end{proof}
\begin{proof}[Proof of the sufficient direction of Theorem A]
By Proposition~\ref{rank}, each element of $B$ has a $p$-power in $R$. Lemma~\ref{wn} gives $B=R$ by iteration, and Proposition~\ref{eval} gives the retraction.
\end{proof}

\begin{corollary}\label{fpure}
If $R$ is one-dimensional Gorenstein of characteristic $p>0$, then weak normality is equivalent to $F$-purity. In particular in characteristic two, $\K$ is also equivalent to $F$-purity. No finiteness condition on Frobenius is required.
\end{corollary}
\begin{proof}
Write $F_*M$ for restriction of scalars along Frobenius. By local duality for a Gorenstein ring \cite{BH}, the injective hull $E$ of $k$ identifies with $H^1_{\m}(R)=Q/R$. Localizing the twisted module gives $Q\otimes_R F_*R\simeq F_*Q$, via $v\otimes F_*a\mapsto F_*(v^p a)$. Taking cokernels of the maps from $F_*R$ therefore identifies $E\otimes_R F_*R$ with $F_*(Q/R)$. Under this identification the map obtained by tensoring Frobenius with $E$ sends $v+R$ to $F_*(v^p+R)$. Its injectivity is precisely root closure.

For completeness, a map $i:R\to M$ is pure if $E\otimes i$ is injective. Indeed the latter injection splits by injectivity of $E$. Adjunction sends a left inverse to an $R$-linear map $M\to\operatorname{End}_R(E)=\widehat R$ whose composite with $i$ is the completion map. After tensoring with any module the latter is still injective by faithful flatness. This proves purity. The converse follows from its definition. The result follows from Lemma~\ref{wn} and Theorem A.
\end{proof}
